\documentclass[a4paper,12pt]{article}
\usepackage[utf8]{inputenc}
\usepackage[T1]{fontenc}
\usepackage[ngerman]{babel}
\usepackage{amsmath}
\usepackage{amssymb}
\usepackage{enumitem}
\usepackage{geometry}
\geometry{a4paper, margin=2.5cm}

\title{Einsendeaufgaben -- Lektion 2\\[0.5em]
\large Modul 61111: Mathematische Grundlagen}
\author{}
\date{}

\begin{document}
\maketitle

\section*{Aufgabe 2.2}

Gauß-Elimination auf der erweiterten Koeffizientenmatrix:

\[
\left(\begin{array}{ccc|c}
2 & 4 & 2 & 12t \\
2 & 12 & 7 & 12t+7 \\
1 & 10 & 6 & 7t+8
\end{array}\right)
\xrightarrow{Z_{2,1}(-1),\; Z_{3,1}\!\left(-\tfrac{1}{2}\right)}
\left(\begin{array}{ccc|c}
2 & 4 & 2 & 12t \\
0 & 8 & 5 & 7 \\
0 & 8 & 5 & t+8
\end{array}\right)
\xrightarrow{Z_{3,2}(-1)}
\left(\begin{array}{ccc|c}
2 & 4 & 2 & 12t \\
0 & 8 & 5 & 7 \\
0 & 0 & 0 & t+1
\end{array}\right)
\]

Damit das Gleichungssystem lösbar ist, muss $t + 1 = 0$, also $t = -1$.

\bigskip
\noindent\textbf{Dann sei $t = -1$:}

\[
\left(\begin{array}{ccc|c}
2 & 4 & 2 & -12 \\
0 & 8 & 5 & 7 \\
0 & 0 & 0 & 0
\end{array}\right)
\xrightarrow{Z_1\!\left(\tfrac{1}{2}\right),\; Z_2\!\left(\tfrac{1}{8}\right)}
\left(\begin{array}{ccc|c}
1 & 2 & 1 & -6 \\
0 & 1 & \tfrac{5}{8} & \tfrac{7}{8} \\
0 & 0 & 0 & 0
\end{array}\right)
\xrightarrow{Z_{1,2}(-2)}
\left(\begin{array}{ccc|c}
1 & 0 & -\tfrac{1}{4} & -\tfrac{31}{4} \\
0 & 1 & \tfrac{5}{8} & \tfrac{7}{8} \\
0 & 0 & 0 & 0
\end{array}\right)
\]

Eine partikuläre Lösung (mit $x_3 = 0$):
\[
d = \begin{pmatrix}-\tfrac{31}{4}\\\tfrac{7}{8}\\0\end{pmatrix}.
\]

\noindent\textbf{Lösung der homogenen Gleichung} ($x_3$ frei):
\[
\left(\begin{array}{ccc}
1 & 0 & -\tfrac{1}{4} \\
0 & 1 & \tfrac{5}{8} \\
0 & 0 & 0
\end{array}\right)
\implies
\mathcal{U} = \left\{
a_1 \begin{pmatrix}-\tfrac{1}{4}\\\tfrac{5}{8}\\-1\end{pmatrix}
\;\middle|\; a_1 \in \mathbb{R}
\right\}.
\]

Die vollständige Lösungsmenge ist:
\[
\mathcal{L} = d + \mathcal{U}
= \left\{
\begin{pmatrix}-\tfrac{31}{4}\\\tfrac{7}{8}\\0\end{pmatrix}
+ a_1 \begin{pmatrix}-\tfrac{1}{4}\\\tfrac{5}{8}\\-1\end{pmatrix}
\;\middle|\; a_1 \in \mathbb{R}
\right\}.
\]

\end{document}
